\clearpage\SourceLeaf{520}
\renewcommand{\BookRunningTitle}{Sectio octava.}
\pdfbookmark[0]{Sectio octava. Quarundam disquisitionum ad circuli sectionem pertinentium uberior consideratio.}{sectio-octava}
\begin{center}
\textbf{SECTIO OCTAVA.}\\[1em]
QUARUNDAM DISQUISITIONUM AD CIRCULI SECTIONEM PERTINENTIUM UBERIOR CONSIDERATIO.
\end{center}
\BookRule{.18\textwidth}
\par\medskip\Needspace{4\baselineskip}\pdfbookmark[1]{367.}{sectio-octava-art-367}\centerline{\textbf{367.}}\nobreak\smallskip
Quae in posteriore Sectionis septimae parte inde ab art $355$ tradidimus, gravia utique specimina exhibent de magna theoriae sectionis circuli fertilitate, nec non de nexu miro, qui hanc disciplinam cum variis disquisitionibus \textit{arithmeticis} iungit. Illic vero, spatii temporisque angustia nimis coarctati, leviter tantum huncce campum stringere potuimus, qui quo ulterius in eo progredimur, eo largiore messe conatus nostros remuneratur. Propositum itaque nobis est, unam alteramve quaestionum ibi inceptarum hic denuo resumere copiosiusque pertractare: certoque lectores non sine magna admiratione plurium problematum arithmeticorum, quae toto hinc coelo dissita esse quisque expectavisset, solutionem huic fundamento inniti videbunt.
\par\medskip\Needspace{4\baselineskip}\pdfbookmark[1]{368.}{sectio-octava-art-368}\centerline{\textbf{368.}}\nobreak\smallskip
Argumentum fertilissimum suppeditat disquisitio in art. $356$ inchoata, ubi complexu radicum aequationis $x^n-1=0$ (unitate exclusa) in duas classes discerpto, aggregatum in utraque classe definire docuimus, quae scilicet prodierunt $=-\frac12+\frac12\surd n$ et $-\frac12-\frac12\surd n$ pro casu ubi $n$ est formae $4n+1$, aut $=-\frac12+\frac12\surd -n$ et $-\frac12-\frac12\surd -n$ pro casu ubi $n$ est formae $4n+3$. Attamen illic non solum limitationem ad casum ubi $n$ est \textit{numerus primus} nobis \SourceLeaf{521}imposueramus, sed etiam, quod multo adhuc gravioris erat momenti, \textit{signum} quantitatis radicalis indefinitum reliquimus, seu potius hanc determinationem paucis addigitatam demonstratione solida fulcire negleximus. Hos itaque defectus ante omnia supplere oportebit.
\par\medskip\Needspace{4\baselineskip}\pdfbookmark[1]{369.}{sectio-octava-art-369}\centerline{\textbf{.369.}}\nobreak\smallskip
Jam sit itaque $n$ numerus integer positivus quicunque, $R$ radix aequationis $x^n-1\pdot=0$ talis, cuius nulla potestas inferior quam $n^{\mathrm{ta}}$ unitati aequalis fiat (V. art. $359$, II.), designemusque per $[\lambda]$, ut in Sect. VII potestatem $R^\lambda$, ita ut $[0]=1$, $[1]$, $[2]$, $[3]\dotsfour[n-1]$ omnes radices aequationis $x^n-1=0$ exhibeant. Porro denotemus aggregatum
\[
{}[0]+[1]+[4]+[9]+\dotsfour+[(n-1)^2]\quad\text{per }\Sigma[\mathfrak{Q}]
\]
et generalius
\[
{}[0]+[\lambda]+[4\lambda]+[9\lambda]\dotsfour+[\lambda(n-1)^2]\quad\text{per }\Sigma[\mathfrak{Q}\lambda]
\]
ita ut $\mathfrak{Q}$ indefinite quadrata numerorum $0,1,2,3\dotsfour n-1$ indicet. Patet igitur, sicut generaliter est $[\lambda]=[\mu]$, si $\lambda$, $\mu$ sunt integri quicunque (positivi seu negativi) secundum $n$ congrui, ita etiam fore $\Sigma[\mathfrak{Q}\lambda]=\Sigma[\mathfrak{Q}\mu]$, si $\lambda\equiv\mu$. His ita praeparatis habemus sequens.
\par\medskip\Needspace{4\baselineskip}\pdfbookmark[1]{370.}{sectio-octava-art-370}\centerline{\textbf{370.}}\nobreak\smallskip
\textsc{Problema.} \textit{Productum e duobus aggregatis $\Sigma[\mathfrak{Q}]$ et $\Sigma[-\mathfrak{Q}]$ assignare.}
\par\noindent\textit{Solutio.} Quum sit $nn\equiv0$, $(n+1)^2\equiv1$, $(n+2)^2\equiv4$ etc. $(\mathrm{mod}.n)$, facile patet fieri $\Sigma[\mathfrak{Q}]$
\[
\begin{aligned}
&=[1]+[4]+[9]+[16]\dotsfive+[nn]\\
&=[4]+[9]+[16]+[25]\dotsfive+[(n+1)]^2\\
&=[9]+[16]+[25]+[36]\dotsfive+[(n+2)^2]
\end{aligned}
\]
\centerline{etc.\quad aut generaliter}
\[
{}=[kk]+[(k+1)^2]+[(k+2)^2]+[(k+3)^2]\dotsfour+[(n+k-1)^2]
\]
Hinc $[-kk]\times\Sigma[\mathfrak{Q}]$
\[
{}=[0]+[2k+1]+[4k+4]+[6k+9]\dotsfive+[(n-1)^2+2(n-1)k]
\]
